% Options for packages loaded elsewhere
\PassOptionsToPackage{unicode}{hyperref}
\PassOptionsToPackage{hyphens}{url}
\documentclass[
]{article}
\usepackage{xcolor}
\usepackage{amsmath,amssymb}
\setcounter{secnumdepth}{-\maxdimen} % remove section numbering
\usepackage{iftex}
\ifPDFTeX
  \usepackage[T1]{fontenc}
  \usepackage[utf8]{inputenc}
  \usepackage{textcomp} % provide euro and other symbols
\else % if luatex or xetex
  \usepackage{unicode-math} % this also loads fontspec
  \defaultfontfeatures{Scale=MatchLowercase}
  \defaultfontfeatures[\rmfamily]{Ligatures=TeX,Scale=1}
\fi
\usepackage{lmodern}
\ifPDFTeX\else
  % xetex/luatex font selection
\fi
% Use upquote if available, for straight quotes in verbatim environments
\IfFileExists{upquote.sty}{\usepackage{upquote}}{}
\IfFileExists{microtype.sty}{% use microtype if available
  \usepackage[]{microtype}
  \UseMicrotypeSet[protrusion]{basicmath} % disable protrusion for tt fonts
}{}
\makeatletter
\@ifundefined{KOMAClassName}{% if non-KOMA class
  \IfFileExists{parskip.sty}{%
    \usepackage{parskip}
  }{% else
    \setlength{\parindent}{0pt}
    \setlength{\parskip}{6pt plus 2pt minus 1pt}}
}{% if KOMA class
  \KOMAoptions{parskip=half}}
\makeatother
\setlength{\emergencystretch}{3em} % prevent overfull lines
\providecommand{\tightlist}{%
  \setlength{\itemsep}{0pt}\setlength{\parskip}{0pt}}
\usepackage{bookmark}
\IfFileExists{xurl.sty}{\usepackage{xurl}}{} % add URL line breaks if available
\urlstyle{same}
\hypersetup{
  hidelinks,
  pdfcreator={LaTeX via pandoc}}

\author{}
\date{}

\begin{document}

\section{Proj-Read-Write-Features
Terminology}\label{proj-read-write-features-terminology}

This a running doc of the terminology convention that I would like to
use and some of the mathematical theory behind things.

\subsection{1. Terminology}\label{terminology}

\begin{itemize}
\tightlist
\item
  \(L\) is the set of layers
\item
  \(J\) is the set of attention heads
\item
  \(d_{\mathrm{model}}\) is residual-stream width.
\item
  \(d_{\mathrm{head}}\) is per-head width.
\item
  \(a_{\ell, A} \in \mathbb{R}^{d_{\mathrm{model}}}\) is the activation
  at layer \(\ell\) for task A
\item
  \(h\) is a head
\item
  \(\mathcal{T}\) is the task universe (antonym, synonym,
  English--French, country--capital, \ldots) and \(A \in \mathcal{T}\)
  is a task.
\item
  \(p_A^j\) is the \(j\)-th prompt for task \(A\)
\item
  \(\mathcal{P}_A = \{p_A^j\}_j\) is the prompt set for \(A\) (varying
  in ICL context length, unless explicitly mentioned otherwise)
\item
  \(h(p_A^j) \in \mathbb{R}^{d_{\mathrm{model}}}\) is the head
  activation at the final token position (final cue token) on prompt
  \(p_A^j\).
\item
  The head vector \(h_A = \frac{1}{|\mathcal{P}_A|} \sum_j h(p_A^j)\) is
  the average head activation for task A.
\item
  \(v_A = \sum_{h \in H} h_A\) is the function vector for task \(A\)
  where \(H\) is the selected subset of heads for our defintion of
  function vectors.
\item
  \(v^j_A = \sum_{h \in H} h(p_A^j)\) is the per prompt function vector
  for prompt \(j\) on task \(A\).
\item
  \(z^t_\ell\) is the residual stream at layer \(\ell\) at token \(t\).
\end{itemize}

\subsection{2. Read Directions}\label{read-directions}

If we treat the span of \(\{v_A\}_{A \in \mathcal{T}}\) as the task
imitation space, then we want to get the task identification space. This
is some causal subspace (ideally ocurring at earlier tokens at earlier
layers) which is feeds into the model machinery to output the task
imitation space. To do this, we can decompose each \(v_A\) into it's
head vectors \(\sum_{h \in H} h_A\). It is then possible to compute the
\emph{read direction} (previously called \emph{d\_payload}) for each
\(h_A\). The idea is that the read direction, \(r^h_A\), is the unit
vector direction which maximises the OV circuit \(W_O W_V\) to output a
vector in the direction of \(h_A\).

\textbf{Definition (Read direction).} Let
\(W_O W_V \in \mathbb{R}^{d_{\mathrm{model}} \times d_{\mathrm{model}}}\)
be the OV circuit of head \(h\). The \emph{read direction}
\(r^h_A \in \mathbb{R}^{d_{\mathrm{model}}}\) for task \(A\) is the
input direction to \(W_O W_V\) whose image is maximally aligned with the
head vector, subject to having no projection along the kernel of
\(W_O W_V\):

\[
r^h_A \;=\; \operatorname*{arg\,max}_{\substack{z \in \mathbb{R}^{d_{\mathrm{model}}},\; \|z\| = 1 \\ z \,\perp\, \ker(W_O W_V)}} \cos\!\left(W_O W_V\, z,\; h_A\right). \tag{1}
\]

The kernel constraint is there to make the argmax unique and well
defined. The OV circuit is low rank since \(W_O\) and \(W_V\) are
rectangular, so \(W_O W_V\) has rank at most \(d_{\mathrm{head}}\) and
its kernel has dimension at least
\(d_{\mathrm{model}} - d_{\mathrm{head}}\) (\(4096 - 256 = 3840\) for
GPT-J). Without the constraint adding any kernel vector \(u\) to a
maximiser \(z^\ast\) leaves the output unchanged
(\(W_O W_V (z^\ast + u) = W_O W_V z^\ast\)), but I would not say that
\(u\) is part of the task identification subspace since it has no causal
downstream impact if it gets ignored by \(M\).

The maximum cosine of \(1\) is attained by the normalised truncated
pseudo inverse (where you do SVD on \(W_O W_V\) and remove all singular
values of 0 and then take the inverse).

\[
r^h_A = \frac{(W_O W_V)^{+}\, h_A}{\|(W_O W_V)^{+}\, h_A\|}. \tag{2}
\]

In practice, \((W_O W_V)^{+}\) scales components by \(1/\sigma_i\), so
near-zero singular values get amplified. (I believe this is the point
that Dan was making earlier but I misunderstood due to the below)

\subsubsection{Previous
misunderstanding}\label{previous-misunderstanding}

\begin{enumerate}
\def\labelenumi{\arabic{enumi})}
\tightlist
\item
  Previously, I then took this definition and conflated maximising
  cosine similarity with the dot product. This gives the solution
  \(r^h_A = (W_O W_V)^\top h_A\) (unit normalised), since that is the
  unit vector maximising the dot product
  \(\langle W_O W_V\, z, h_A \rangle = \langle z, (W_O W_V)^\top h_A \rangle\).
  I assumed the unit-norm constraint on \(z\) made the two equivalent.
  It does not (the constraint fixes the norm of the input \(z\), but the
  cosine normalises by the norm of the output \(W_O W_V z\), which still
  varies with \(z\)). The dot product rewards output magnitude while the
  cosine is indifferent to it, so the two problems have different
  solutions, and the transpose formula solves only the dot-product one.
\item
  I thought that if the task imitation space is \(m\)-dimensional,
  e.g.~spanned by \(\{v_{j}\}_{j = 1 \colon m}\), then the task
  identification space can be spanned by at most \(m\) vectors
  \(\{(W_O W_V)^{+} v_{j}\}_{j = 1 \colon m}\). But this is only
  inverting the task subspace for a single head when in reality there
  are multiple heads, so you can invert a single direction in the task
  imitation space \(v_j\) as \(\{(W^h_O W^h_V)^{+} v_{j}\}_{h \in H}\)
  which means that you can have a larger dimension task identification
  subspace.
\end{enumerate}

\subsection{3. Per Prompt Read
Direction}\label{per-prompt-read-direction}

One of the issues mentioned in the previous section is that there are
multiple heads invloved in the computation of one function vector. After
some thought, I had the following idea for \emph{per prompt read
directions}, which are basically just the read direction that maximises
the cosine similarity to the per prompt function vectors when passed
through all of the heads that are selected. To formulate this
mathematically, let

\[
M \;=\; \sum_{h \in H} W_O^h W_V^h \;\in\; \mathbb{R}^{d_{\mathrm{model}} \times d_{\mathrm{model}}} \tag{3}
\]

be the summed OV circuit of the selected FV heads \(H\). (Note that
summing across heads treats every head as reading the \emph{same}
residual-stream vector \(z\), even though the heads sit at different
layers)

\textbf{Definition (Per-prompt read direction).} The \emph{per-prompt
read direction} \(r^j_A \in \mathbb{R}^{d_{\mathrm{model}}}\) for prompt
\(p_A^j\) is the unit input direction to \(M\) whose image is maximally
aligned with the per-prompt function vector \(v^j_A\), subject to having
no projection along the kernel of \(M\):

\[
r^j_A \;=\; \operatorname*{arg\,max}_{\substack{z \in \mathbb{R}^{d_{\mathrm{model}}},\; \|z\| = 1 \\ z \,\perp\, \ker(M)}} \cos\!\left(M z,\; v^j_A\right). \tag{4}
\]

As in the single-head case, the maximum is attained by the normalised
truncated pseudo inverse:

\[
r^j_A \;=\; \frac{M^{+}\, v^j_A}{\|M^{+}\, v^j_A\|}. \tag{5}
\]

Replacing \(v^j_A\) with \(v_A\) gives the corresponding task-level read
direction \(r_A\) under \(M\).

Again, using the no projection in the kernel space ensures uniqueness,
so that each prompt has one read direction.

\textbf{Caveat (near-kernel).} The kernel constraint resolves
multiplicity from \emph{exact} zeros only. A full-rank \(M\) can still
have very near zero singular values: \(M^{+}\) scales by
\(1/\sigma_i\), so the exact solution can be dominated by directions the
circuit barely transmits. There might be some optimisation do to around
this: e.g.~a robust variant could be the \(\tau\)-truncated pseudo
inverse \(M^{+}_\tau\) (invert only \(\sigma_i \ge \tau\); equivalently
\(z \perp\) the right-singular subspace with \(\sigma_i < \tau\)).

\end{document}
